\newpage
% =================================================
\section{Materialized (Views)}
% =================================================
\subsection{Tabela Materialized para Lote}
\begin{entidade}{Lote\_Materializado}
\campo{id}{SERIAL}{PK}
\campo{id\_lote}{INTEGER}{FK, UNIQUE}
\campo{qtd\_frascos\_cadastrados}{INTEGER}{DEFAULT 0, NOT NULL}
\campo{ultimo\_reconciliador}{TIMESTAMP}{NULL}
\end{entidade}
\begin{explicacao}
Corrigido o typo \texttt{SERAIL} $\rightarrow$ \texttt{SERIAL}. \textit{id\_lote} recebeu \texttt{UNIQUE}: esta é uma tabela 1:1 com \textit{Lote} (uma linha por lote), então o par (\textit{id}, \textit{id\_lote}) sem essa restrição permitiria duas linhas materializadas divergentes para o mesmo lote. Mantida \textbf{exclusivamente pela transação síncrona} da \textit{Cloud Function} a cada frasco cadastrado, sem depender de triggers assíncronas. Adicionada a marca d'água \textit{ultimo\_reconciliador} que é preenchida por \textit{jobs} agendados de reconciliação para bloquear o duplo cômputo de eventos atrasados (drift silenciado, ver seção \textit{Tecnologia e Relatórios}).
\end{explicacao}

\subsection{Tabela Resumo de Almoxarifado - Diário}
\begin{entidade}{Resumo\_Almoxarifado\_Diario}
\campo{id}{SERIAL}{PK}
\campo{data}{DATE}{NOT NULL}
\campo{id\_almoxarifado}{INTEGER}{FK, NOT NULL}
\campo{qtd\_frascos\_cadastrados\_fechados\_durante\_o\_dia}{INTEGER}{DEFAULT 0, NOT NULL}
\campo{qtd\_frascos\_cadastrados\_abertos\_durante\_o\_dia}{INTEGER}{DEFAULT 0, NOT NULL}
\campo{qtd\_frascos\_emprestados\_durante\_o\_dia}{INTEGER}{DEFAULT 0, NOT NULL}
\campo{qtd\_frascos\_devolvidos\_durante\_o\_dia}{INTEGER}{DEFAULT 0, NOT NULL}
\campo{volume\_utilizado\_no\_dia\_ml}{NUMERIC(12,2)}{DEFAULT 0, NOT NULL}
\campo{volume\_evaporado\_no\_dia\_ml}{NUMERIC(12,2)}{DEFAULT 0, NOT NULL}
\campo{massa\_utilizada\_no\_dia\_g}{NUMERIC(12,2)}{DEFAULT 0, NOT NULL}
\campo{massa\_evaporada\_no\_dia\_g}{NUMERIC(12,2)}{DEFAULT 0, NOT NULL}
\campo{calculado\_em}{TIMESTAMP}{NOT NULL -- serverTimestamp da execução que produziu a projeção atual}
\campo{versao\_calculo}{TEXT}{NOT NULL -- identificador da versão da regra/algoritmo de materialização}
\campo{created\_at}{TIMESTAMP}{NOT NULL}
\campo{updated\_at}{TIMESTAMP}{NOT NULL}
\end{entidade}
\begin{regra}
\textit{UNIQUE(data, id\_almoxarifado)}. No Firestore, o documento usa ID determinístico \texttt{YYYY-MM-DD\_\_ALMOXARIFADO\_ID}. Reprocessar a mesma chave diária substitui a projeção existente; nunca cria um novo documento.
\end{regra}
\newpage
\begin{explicacao}[Separação Física Rigorosa: Massa ($g$) vs Volume ($ml$)]
Esta tabela consolida exclusivamente fatos FLOW do dia civil, sem posição acumulada de estoque.
Para eliminar distorções métricas em sólidos com densidade desconhecida ou nula, as colunas agregadas de consumo e evaporação são formalmente segregadas nas duas dimensões físicas fundamentais: reagentes líquidos acumulam em \textit{volume\_...\_ml} e reagentes sólidos acumulam em \textit{massa\_...\_g}. Nunca se tenta converter gravimetria de sólidos para volume sem densidade aferida.

O campo legado que misturava unidades foi removido por redundância; consumo usa exclusivamente as métricas dimensionadas em g e mL.
\end{explicacao}

\begin{explicacao}[Projeção FLOW e Unicidade Determinística]
\textit{Resumo\_Almoxarifado\_Diario} é uma \textbf{projeção derivada, descartável, recalculável e reconstruível} a partir das fontes históricas canônicas (\textit{Historico\_Frasco\_Reagente}, \textit{Emprestimo\_Reagente} e eventos de domínio; ver Catálogo de Métricas). \textbf{Não é fonte primária de verdade.} Não se admite correção manual incremental de totais como autoridade: qualquer divergência deve ser resolvida por reprocessamento absoluto a partir das fontes.

\textbf{Unicidade e ID determinístico.} A chave lógica diária é \textit{UNIQUE(data, id\_almoxarifado)}. No Firestore o ID do documento é determinístico no formato \texttt{YYYY-MM-DD\_\_ALMOXARIFADO\_ID}. Reprocessar a mesma data/almoxarifado deve \textbf{substituir} (sobrescrever) a projeção existente, jamais criar um novo documento por reprocessamento.

\textbf{Metadados de cálculo.} \textit{data} identifica o dia civil agregado; \textit{calculado\_em} é o \texttt{serverTimestamp} da execução que produziu a projeção atual; \textit{versao\_calculo} identifica a versão da regra/algoritmo de materialização. \textit{calculado\_em} \textbf{não é a data do fato}: não se confunde com \textit{data} nem com o timestamp físico dos eventos agregados.

\textbf{Janela FLOW.} O dia civil D em America/Sao\_Paulo usa o intervalo semiaberto do início de D até o início de D+1. Somente fatos efetivos pertencentes a D são agregados. Não há snapshot diário de estoque, leitura de estado corrente para fabricar passado, dependência de D-1 ou replay até hoje. Correções fechadas reprocessam somente as datas diretamente afetadas (origem e destino, se houver mudança de atribuição).

\textbf{Correção de fato histórico.} Uma correção administrativa que altera um fato histórico usado na agregação deve identificar a \textit{DATA\_ALVO} do fato corrigido e disparar o reprocessamento do resumo de \textbf{cada data histórica afetada}. O timestamp da correção permanece o da correção; o timestamp físico original não é reescrito.
\end{explicacao}
\newpage
\subsection{Tabela Resumo de Reagente/Solvente - Diário}
\begin{entidade}{Resumo\_Reagente\_Diario}
\campo{id}{SERIAL}{PK}
\campo{data}{DATE}{NOT NULL}
\campo{id\_resumo\_reagente}{INTEGER}{FK, NOT NULL}
\campo{id\_almoxarifado}{INTEGER}{FK, NOT NULL}
\campo{qtd\_frascos\_descartados\_dia}{INTEGER}{DEFAULT 0, NOT NULL}
\campo{qtd\_frascos\_extraviados\_dia}{INTEGER}{DEFAULT 0, NOT NULL}
\campo{volume\_utilizado\_no\_dia\_ml}{NUMERIC(12,2)}{DEFAULT 0, NOT NULL}
\campo{volume\_evaporado\_no\_dia\_ml}{NUMERIC(12,2)}{DEFAULT 0, NOT NULL}
\campo{massa\_utilizada\_no\_dia\_g}{NUMERIC(12,2)}{DEFAULT 0, NOT NULL}
\campo{massa\_evaporada\_no\_dia\_g}{NUMERIC(12,2)}{DEFAULT 0, NOT NULL}
\campo{calculado\_em}{TIMESTAMP}{NOT NULL -- serverTimestamp da execução que produziu a projeção atual}
\campo{versao\_calculo}{TEXT}{NOT NULL -- identificador da versão da regra/algoritmo de materialização}
\campo{created\_at}{TIMESTAMP}{NOT NULL}
\campo{updated\_at}{TIMESTAMP}{NOT NULL}
\end{entidade}
\begin{regra}
\textit{UNIQUE(data, id\_almoxarifado, id\_resumo\_reagente)}. No Firestore, o documento usa ID determinístico \texttt{YYYY-MM-DD\_\_ALMOXARIFADO\_ID\_\_RESUMO\_REAGENTE\_ID}. Reprocessar a mesma chave diária substitui a projeção existente; nunca cria um novo documento.
\end{regra}

\begin{explicacao}[Contratos do Estoque Atual (fora dos resumos diários)]
As regras a seguir pertencem exclusivamente às agregações atuais protegidas no backend sobre Frasco\_Reagente, não aos resumos FLOW:
\begin{itemize}
    \item \textbf{Frascos Extraviados:} Frascos com estado \texttt{EXTRAVIADO} \textbf{nunca} são contabilizados no saldo disponível, seja na quantidade de frascos ou na massa/volume total.
    \item \textbf{Conteúdo Nominal vs. Saldo:} O \textit{conteudo\_nominal} (capacidade original do fabricante) \textbf{não representa} o saldo atual do frasco e não deve ser somado no lugar de saldos medidos.
    \item \textbf{Saldos Desconhecidos:} Se o saldo de um frasco é desconhecido (ex.: frasco aberto sem \textit{peso\_atual} preciso registrado), ele não deve ser somado como zero (\texttt{0}) nos agregados totais de volume ou massa, pois isso distorceria a realidade física. Esses frascos devem ser segregados na camada de agregação, distinguindo claramente a ``quantidade conhecida'' da ``quantidade com saldo desconhecido''.
\end{itemize}
\end{explicacao}

\begin{explicacao}[Projeção FLOW e Unicidade Determinística]
\textit{Resumo\_Reagente\_Diario} é uma \textbf{projeção derivada, descartável, recalculável e reconstruível} a partir das fontes históricas canônicas (\textit{Historico\_Frasco\_Reagente}, \textit{Emprestimo\_Reagente} e eventos de domínio). \textbf{Não é fonte primária de verdade.} Não se admite correção manual incremental de totais como autoridade: divergências são resolvidas por reprocessamento absoluto a partir das fontes.

\textbf{Unicidade e ID determinístico.} A chave lógica diária é \textit{UNIQUE(data, id\_almoxarifado, id\_resumo\_reagente)}. No Firestore o ID do documento é determinístico no formato \texttt{YYYY-MM-DD\_\_ALMOXARIFADO\_ID\_\_RESUMO\_REAGENTE\_ID}. Reprocessar a mesma chave deve \textbf{substituir} (sobrescrever) a projeção existente, jamais criar um novo documento por reprocessamento.

\textbf{Metadados de cálculo.} \textit{data} identifica o dia civil agregado; \textit{calculado\_em} é o \texttt{serverTimestamp} da execução que produziu a projeção atual; \textit{versao\_calculo} identifica a versão da regra/algoritmo de materialização. \textit{calculado\_em} \textbf{não é a data do fato}.

\textbf{Janela FLOW.} O dia civil D em America/Sao\_Paulo usa o intervalo semiaberto do início de D até o início de D+1. Somente fatos efetivos pertencentes a D são agregados. Não há snapshot diário de estoque, leitura de estado corrente para fabricar passado, dependência de D-1 ou replay até hoje. Correções fechadas reprocessam somente as datas diretamente afetadas (origem e destino, se houver mudança de atribuição).

\textbf{Correção de fato histórico.} Uma correção administrativa que altera um fato histórico agregado deve identificar a \textit{DATA\_ALVO} do fato corrigido e reprocessar o resumo de \textbf{cada data histórica afetada}, sem reescrever o timestamp físico original.
\end{explicacao}

\subsection{Tabela Resumo de Bens Patrimoniais - Diário}
\begin{entidade}{Resumo\_Bem\_Patrimonial\_Diario}
\campo{id}{SERIAL}{PK}
\campo{data}{DATE}{NOT NULL}
\campo{id\_local}{INTEGER}{FK, NULL}
\campo{qtd\_bens\_ativos\_no\_fim\_do\_dia}{INTEGER}{DEFAULT 0, NOT NULL}
\campo{qtd\_bens\_inservivel\_no\_fim\_do\_dia}{INTEGER}{DEFAULT 0, NOT NULL}
\campo{qtd\_bens\_dado\_baixa\_no\_fim\_do\_dia}{INTEGER}{DEFAULT 0, NOT NULL}
\campo{qtd\_bens\_cadastrados\_durante\_o\_dia}{INTEGER}{DEFAULT 0, NOT NULL}
\campo{qtd\_bens\_marcados\_inservivel\_durante\_o\_dia}{INTEGER}{DEFAULT 0, NOT NULL}
\campo{qtd\_bens\_dados\_baixa\_durante\_o\_dia}{INTEGER}{DEFAULT 0, NOT NULL}
\campo{qtd\_requisicoes\_abertas\_durante\_o\_dia}{INTEGER}{DEFAULT 0, NOT NULL}
\campo{qtd\_requisicoes\_aprovadas\_durante\_o\_dia}{INTEGER}{DEFAULT 0, NOT NULL}
\campo{qtd\_requisicoes\_rejeitadas\_durante\_o\_dia}{INTEGER}{DEFAULT 0, NOT NULL}
\campo{created\_at}{TIMESTAMP}{NOT NULL}
\campo{updated\_at}{TIMESTAMP}{NOT NULL}
\end{entidade}

\subsection{Tabela de Atividade Mensal - Gestor de Almoxarifado}
\begin{entidade}{Atividade\_Gestor\_Almoxarifado\_Mensal}
\campo{id}{SERIAL}{PK}
\campo{id\_gestor}{INTEGER}{FK, NOT NULL}
\campo{ano}{INTEGER}{NOT NULL}
\campo{mes}{INTEGER}{NOT NULL}
\campo{qtd\_reagentes\_adicionados}{INTEGER}{DEFAULT 0, NOT NULL}
\campo{qtd\_frascos\_adicionados}{INTEGER}{DEFAULT 0, NOT NULL}
\campo{qtd\_movimentacoes\_registradas}{INTEGER}{DEFAULT 0, NOT NULL}
\campo{created\_at}{TIMESTAMP}{NOT NULL}
\campo{updated\_at}{TIMESTAMP}{NOT NULL}
\end{entidade}
\begin{regra}
\textit{UNIQUE(id\_gestor, ano, mes)}.
\end{regra}

\subsection{Tabela de Atividade Mensal - Gestor de Bens Patrimoniais}
\begin{entidade}{Atividade\_Gestor\_Bens\_Patrimoniais\_Mensal}
\campo{id}{SERIAL}{PK}
\campo{id\_gestor}{INTEGER}{FK, NOT NULL}
\campo{ano}{INTEGER}{NOT NULL}
\campo{mes}{INTEGER}{NOT NULL}
\campo{qtd\_bens\_adicionados}{INTEGER}{DEFAULT 0, NOT NULL}
\campo{qtd\_bens\_editados}{INTEGER}{DEFAULT 0, NOT NULL}
\campo{qtd\_requisicoes\_aprovadas}{INTEGER}{DEFAULT 0, NOT NULL}
\campo{qtd\_requisicoes\_rejeitadas}{INTEGER}{DEFAULT 0, NOT NULL}
\campo{created\_at}{TIMESTAMP}{NOT NULL}
\campo{updated\_at}{TIMESTAMP}{NOT NULL}
\end{entidade}
\begin{regra}
\textit{UNIQUE(id\_gestor, ano, mes)}.
\end{regra}

\subsection{Catálogo de Métricas Disponíveis por Domínio}
\begin{itemize}
    \item \textbf{\textit{Historico\_Frasco\_Reagente}} (por \textit{tipo}): frascos cadastrados, saídas, entradas, esvaziados, quebrados, descartados, vencidos, ajustes de inventário e evaporação -- por dia $\rightarrow$ alimenta \textit{Resumo\_Almoxarifado\_Diario} e \textit{Resumo\_Reagente\_Diario}; por gestor/mês $\rightarrow$ alimenta \textit{Atividade\_Gestor\_Almoxarifado\_Mensal}; contagem "agora" $\rightarrow$ \texttt{count()} server-side via backend protegido sobre \textit{Frasco\_Reagente} filtrado por \textit{vencido = true}.
    \item \textbf{\textit{Emprestimo\_Reagente}}: empréstimos abertos/devolvidos/atrasados por dia, volume e massa de consumo validado por dia $\rightarrow$ alimenta os resumos diários de reagente.
    \item \textbf{\textit{Historico\_Bem\_Patrimonial} / \textit{Alteracao\_Bem\_Patrimonial}}: bens cadastrados, editados e baixados por dia $\rightarrow$ alimenta \textit{Resumo\_Bem\_Patrimonial\_Diario}.
    \item \textbf{\textit{Requisicao\_Edicao\_Bem\_Patrimonial} / \textit{Requisicao\_Adicao\_Bem\_Patrimonial}}: requisições abertas/aprovadas/rejeitadas por dia $\rightarrow$ alimenta \textit{Resumo\_Bem\_Patrimonial\_Diario}.
    \item \textbf{\textit{Registro\_de\_Auditoria}}: log genérico entre todas as entidades.
\end{itemize}

\subsection{Estoque atual: autoridade canônica e agregação protegida}
\label{sec:estoque-atual-cache}
\textit{Frasco\_Reagente} é a autoridade do estado corrente. Não existe materialized view persistente de estoque atual. O backend executa \texttt{count()} e \texttt{sum()} sobre esse documento canônico; o navegador não executa agregações Firestore nem baixa frascos para contá-los. O contador de vínculos de lote tem finalidade cadastral distinta e não substitui estoque atual.

Para permitir \texttt{sum()}, a projeção Firestore do próprio frasco possui \texttt{saldo\_aferido\_g} e \texttt{saldo\_aferido\_ml}, derivados server-owned. O backend calcula e atualiza esses campos na mesma transação das alterações de peso, tara ou conhecimento: saldo desconhecido ou metrologicamente inconsistente produz NULL; vazio explicitamente confirmado produz zero; caso consistente usa peso menos tara e, para líquidos, divisão pela densidade válida. A origem da tara permanece explícita. Sólidos têm campo mL NULL. As somas de estoque utilizável filtram aptidão operacional (estado FECHADO/ABERTO, disponível, sem quarentena e respeitando validade/autorização); contagens totais, emprestados e desconhecidos usam filtros próprios identificados. Nenhuma quantidade desconhecida é apresentada como zero conhecido.

\paragraph{Contrato de consulta e cache.}
A callable de dashboard exige App Check, Firebase Auth, RBAC e escopo antes de aplicar limite atômico de \textbf{5 solicitações/minuto por UID}, inclusive em cache hit. Usar janela deslizante de 60 segundos, relógio do servidor e registro interno por UID; IP não é identidade primária. Rejeição ocorre antes de executar agregações. O cache compartilhado nunca dispensa autorização individual.

\texttt{Sistema\_Cache\_Dashboard} é coleção interna, efêmera, derivada e descartável. Chaves permitidas: \texttt{ESTOQUE\_\_ALMOX} e \texttt{ESTOQUE\_\_ALMOX\_\_RESUMO}, com IDs codificados deterministicamente sem colisão. O endpoint aceita somente esses escopos e filtros enumerados; qualquer filtro que altere resultado entra na chave normalizada, sem parâmetros livres que criem cardinalidade ilimitada. Campos: \texttt{escopo}, \texttt{resultado}, \texttt{calculado\_em}, \texttt{versao\_calculo}, \texttt{geracao} e \texttt{valido}. Resultado é nullable durante invalidação. A geração é token de invalidação, não métrica de estoque.

Validade exige idade estritamente menor que 30 segundos, versão compatível e ausência de invalidação posterior. A marca temporal é capturada no início do cálculo (limite conservador da idade dos dados). Remoção física por TTL do Firestore é apenas limpeza, nunca mecanismo preciso de validade. Não existe scheduler de 30 segundos: cache miss recalcula sob demanda; oito horas sem acesso geram zero recálculos.

\paragraph{Invalidação e concorrência.}
Toda mutação que altere contribuição atual atualiza, atomicamente com o frasco, as gerações das chaves do almoxarifado e do par almoxarifado/resumo afetados, marca \texttt{valido=false} e limpa o resultado. Não calcula agregações nessa transação. A invalidação cria o registro de geração se necessário, evitando corrida com primeiro preenchimento. Transferência ou alteração de associação, somente quando existir contrato suportado, invalida origem e destino. A publicação do cache relê a geração e só publica se ela não mudou desde antes das consultas; descarta cálculo obsoleto e tenta novamente de forma limitada. Agregações concorrentes ocasionais são aceitáveis; não introduzir Redis, locks distribuídos ou infraestrutura dedicada. Se a geração mudar repetidamente, responder indisponibilidade temporária em vez de publicar dado invalidado. A remoção física de um registro faz a comparação falhar; nunca tratar ausência como geração válida capturada.

Invalidam cadastro, abertura, retirada, devolução normal ou anômala, resolução de anomalia, esgotamento, pesagem que muda saldo, evaporação, ajuste metrológico, recalibração, quebra, descarte, extravio, reencontro, entrada/saída de quarentena e decisões de validade/descarte que alterem elegibilidade. Atualizações de validade pelo job também invalidam. Correção exclusivamente histórica de retirante, gestor, motivo ou auditoria e reprocessamento FLOW não invalidam estoque atual nem alteram pesos correntes. O menor escopo seguro é sempre o das chaves efetivamente afetadas.

\begin{lstlisting}[language=TypeScript, title={Consulta lazy e publicação condicionada à geração (pseudocódigo)}]
// App Check aplicado pela callable antes do handler.
async function consultarEstoque(request) {
  const escopo = await autenticarAutorizarENormalizarEscopo(request);
  await consumirLimiteUID(request.auth.uid, 5, 60); // atomico; inclui cache hit
  const ref = chaveCache(escopo);
  for (let tentativa = 0; tentativa < 2; tentativa++) {
    const base = await obterOuCriarGeracao(ref);
    if (cacheValido(base, agoraServidor(), 30)) return base;
    const inicio = agoraServidor();
    const resultado = await agregarCountSumFrascos(escopo);
    const publicado = await db.runTransaction(async tx => {
      const atual = await tx.get(ref);
      if (!atual.exists || atual.data().geracao !== base.geracao) return false;
      if (agoraServidor() - inicio >= 30000) return false;
      tx.set(ref, { escopo, resultado, calculado_em: inicio,
        versao_calculo: "estoque-v1", geracao: base.geracao, valido: true });
      return true;
    });
    if (publicado) return { resultado, calculado_em: inicio };
  }
  throw new HttpsError("unavailable", "Estoque mudou; tente novamente.");
}
\end{lstlisting}

\paragraph{Interface e relatórios.}
A UI mantém cache curto, nunca além da validade restante informada pelo servidor, cancela requisições antigas ao mudar escopo e invalida seu cache após mutação. Exibe instante de cálculo, atualização em andamento e ação Atualizar com tratamento de rate limit; não consulta automaticamente a cada 30 segundos. Dados já exibidos podem ficar desatualizados por mutação de outro usuário, portanto nunca autorizam retirada. Estado atual é fotografia corrente separada. Relatórios por período usam FLOW e fatos históricos; não prometem séries de saldo, fechados ou emprestados ao fim de cada dia. Patrimônio mantém seu contrato próprio de snapshots diários.
